\documentclass{article}

\usepackage[preprint]{neurips_2023}
\usepackage{hyperref}
\usepackage{url}
\usepackage{booktabs}
\usepackage{amsfonts}

\usepackage{nicefrac}
\usepackage{microtype}
\usepackage{xcolor}
\usepackage{graphicx}
\usepackage{algorithm}
\usepackage{algorithmic}
\usepackage{dirtree}

\usepackage[utf8]{inputenc}
\usepackage[T1]{fontenc}
\usepackage{amsmath,amssymb,amsthm}
\usepackage{mathtools}
\usepackage{enumitem}
\usepackage{hyperref}
\usepackage{booktabs}

\usepackage[most]{tcolorbox}
\tcbuselibrary{skins, breakable}
\tcbset{
  aclbox/.style={
    colback=gray!5,
    colframe=black!30,
    boxrule=0.5pt,
    arc=2pt,
    left=2pt,right=2pt,top=2pt,bottom=2pt,
  }
}


\theoremstyle{definition}
\newtheorem{definition}{Definition}[section]
\newtheorem{remark}{Remark}[section]
\newtheorem{example}{Example}[section]

\title{\textsc{Struct Memory-R1}: Empowering Large Language Model Agents with Hierarchical Memory Organization via Reinforcement Learning}

\author{
  Yifan Liu \And Liam Gallagher \And David Courtis \And Jiazhou Liang
}

\begin{document}

\maketitle

%% =====================================================================
%% SECTION 1: PROBLEM DEFINITION
%% =====================================================================
\section{Introduction}
Large Language Models (LLM) have evolved from static text generators into interactive agents that engage in multi-turn dialogue, tool use, and long-horizon reasoning. In these settings, memory plays a central role: agents must retain user preferences, task constraints, and intermediate state in order to respond coherently and consistently over time. Consequently, recent systems increasingly augment LLMs with external memory modules, including vector databases, and interaction histories.\\

Despite this progress, most existing memory-augmented LLMs treat memory access as a static or heuristic process~\citep{lewis2020retrieval}. Memory is typically retrieved via similarity search or rule-based triggers~\citep{khandelwal2020generalization}, even when memory is structured, and its usage is largely decoupled from the model’s reasoning process. As a result, current approaches often suffer from over-retrieval, under-utilization of relevant information, and brittle behavior when memory contains stale or weakly relevant entries. More fundamentally, these systems lack a mechanism to learn when, how, and at what level of structure or abstraction memory should be accessed in a task-dependent manner.\\

Recent work, most notably \citet{jin2025search}, addresses this limitation by treating search as an explicit action within the model’s reasoning process and training a policy over reasoning–search trajectories using relative reinforcement learning. By comparing alternative trajectories for the same query, Search-R1 learns when search is beneficial without relying on absolute reward calibration. Motivated by this paradigm, recent efforts have begun to explore similar ideas for memory usage~\citep{yan2025memory}. However, existing approaches operate over flat or unstructured memory and do not address the challenges of structured memory access.\\

In this work, we observe that memory access is structurally analogous to search, with the key difference that memory operates over an agent’s internal, persistent, and often structured state~\citep{park2023generative}, rather than an external corpus. In contrast to search over unstructured documents, memory access in practical systems involves navigating, selecting, and aggregating over structured representations, such as profiles, task hierarchies, and interaction histories. Both search and memory access involve non-differentiable decisions that trade off informativeness, cost, and downstream performance. This motivates \emph{Structured Memory-R1}, which draws direct inspiration from Search-R1 and applies relative reinforcement learning to structured memory access and manipulation. We treat memory usage as a sequence of explicit, structure-aware actions during generation and optimize these decisions by comparing alternative structured memory–interaction trajectories under relative judgments.\\

\section{Related Work}

\paragraph{Flat and Unstructured Memory.}
Early memory-augmented LLMs retrieve information from flat, unstructured stores via similarity search or rule-based triggers~\citep{lewis2020retrieval, khandelwal2020generalization}.
Agent systems such as Generative Agents~\citep{park2023generative} popularized the episodic memory stream, a flat chronological log retrieved by recency and relevance.
These approaches decouple memory access from the model's reasoning process, leading to over-retrieval and under-utilization of relevant information.
A compounding failure mode is the ``lost in the middle'' phenomenon~\citep{liu2024lost}, wherein serializing the full memory state into the context window causes models to neglect centrally positioned entries---motivating selective, reasoning-driven retrieval over both flat and in-context memory baselines.

\paragraph{Structured Memory Access.}
To address the limitations of flat retrieval, recent work has proposed structured memory representations that expose relational and hierarchical organization to the agent.
\citet{semanticxpath2026} introduce \textsc{SemanticXPath}, a paradigm in which the LLM generates XPath-style queries over a structured memory tree, enabling precise, path-directed access without flat similarity search.
While this approach exploits memory structure at inference time, it relies on heuristic query generation and does not learn an access policy from task feedback.

\paragraph{RL-Based Memory and Search.}
\citet{jin2025search} address the analogous limitation in web search by treating search as an explicit reasoning action and training a policy over reasoning--search trajectories via relative reinforcement learning (GRPO).
Motivated by this paradigm, \citet{yan2025memory} extend the idea to flat agent memory, training a two-agent system (Memory Manager + Answer Agent) to perform ADD/UPDATE/DELETE/NOOP operations via RL.
However, both systems operate over flat or unstructured memory and cannot exploit the hierarchical organization present in structured memory stores.

In this work, we observe that memory access is structurally analogous to search, with the key difference that memory operates over an agent's internal, persistent, and hierarchically organized state~\citep{park2023generative}.
Unlike retrieved web documents, memory entries carry structural provenance---their utility depends not only on semantic relevance but on their position within the memory tree.
\textsc{Struct Memory-R1} addresses this gap by extending the RL-over-actions paradigm to structured memory, training agents to learn \emph{where} in the hierarchy to read and write, and \emph{when} structural reorganization is warranted, entirely through end-to-end reinforcement learning over hierarchical memory trajectories.
\section{Problem Definition}

\paragraph{Conversational Agent.}

We consider an LLM-based agent augmented with an external memory system that persists across interactions:
\begin{equation}
\mathcal{A} = (\pi_\theta,\; \mathcal{X},\; \mathcal{Y},\; \mathcal{M},\; \mathcal{E}),
\end{equation}
\begin{itemize}[nosep]
    \item \textbf{policy} $\pi_\theta$: maps an observation and a memory state to a distribution over action trajectories;
    \item \textbf{input (observation) space}  $\mathcal{X}$: the set of all possible user queries and dialogue the agent may receive;
    \item \textbf{output (response) space} $\mathcal{Y}$ : the set of all possible natural-language responses the agent may produce;
    \item  \textbf{memory space} $\mathcal{M}$: the set of all possible external memory states the agent can read from and write to;
    \item \textbf{external environment interface} $\mathcal{E}$: the set of tools through which the agent can perform actions beyond pure language generation (e.g., a memory retrieval server, a search engine, or a code executor).
\end{itemize}


An \emph{interaction} is a temporally ordered sequence of turns between a user and the agent:
\begin{equation}
\mathcal{I} = (t_1, t_2, \ldots, t_T),
\end{equation}
where $T$ is the total number of turns in the interaction.
At turn $t$ ($1 \le t \le T$), the user provides an input $x_t \in \mathcal{X}$ (a query, instruction, or follow-up). The agent observes $x_t$ together with the current memory state $M_t \in \mathcal{M}$. T

The agent generates an action trajectory as a sequence of \emph{language actions} (reasoning steps, final answers) and \emph{environment actions} (memory queries, tool calls).
The agent emits a final response $y_t \in \mathcal{Y}$ extracted from $\tau_t$. The memory state may be updated: $M_{t+1} = f_{\text{update}}(M_t, \tau_t)$, where $f_{\text{update}}$ encodes any write operations performed during the turn.
The objective is to learn a policy $\pi_\theta$ that controls memory access when and how memory should be queried to improve downstream task performance. 

\paragraph{Memory Structures}

Recent work, notably Search-R1~\citep{jin2025search}  and Memory-R1~\citep{yan2025memory}\footnote{Memory-R1 does not provide a public codebase. We built a publicly available implementation at \url{https://github.com/YifanLiu2/Struct_Memory_R1.git}.}, has shown that RL can effectively train agents to manage flat external memory banks.
A \emph{flat memory state} is an unordered collection of $N$ memory entries:
\begin{equation}
M_t^{\text{flat}} = \{m_1, m_2, \ldots, m_N\}, \quad m_i = (\text{key}_i,\; \text{attributes}_i),
\end{equation}


\section{\textsc{Struct Memory-R1}}

\subsection{Memory States}
Building on Search-R1 and Memory-R1, our work extends reinforcement learning with verifiable rewards from external information retrieval to \emph{structured} memory for long-context conversational AI scenarios.
Specificly, a \emph{structured memory state} at turn $t$ is a rooted tree:
\begin{equation}
M_t = (V_t,\; E_t,\; r),
\end{equation}

$r \in V_t$ is the designated \textbf{root node}.
$V_t$ is a finite set of \textbf{nodes} (memory units). Each node $v \in V_t$ carries:
    \begin{itemize}[nosep]
        \item a \textbf{type label} $\ell(v) \in \mathcal{L}$ (e.g., \textsc{Itinerary}, \textsc{Day}, \textsc{POI}, \textsc{Speaker}, \textsc{Fact});
        \item a set of \textbf{textual attributes} $\text{attr}(v) = \{(k_1, val_1), \ldots, (k_m, val_m)\}$.
    \end{itemize}

 $E_t \subseteq V_t \times V_t$ is a set of directed \textbf{parent--child edges} encoding hierarchical relations such as:
\begin{equation}
\textsc{Itinerary} \;\to\; \textsc{Version} \;\to\; \textsc{Day} \;\to\; \textsc{POI},
\end{equation}


\subsection{Action Space}

The agent's action space is $\mathcal{A} = \mathcal{A}_{\text{lang}} \cup \mathcal{A}_{\text{env}}$, where:
$\mathcal{A}_{\text{lang}}$: \textbf{Reasoning} (\texttt{<think>}) and \textbf{Answer} (\texttt{<answer>}) generation.
And $\mathcal{A}_{\text{env}}$: \textbf{Memory Query} (\texttt{<memory>}$q$\texttt{</memory>}) and \textbf{Memory Update} (\textsc{Add}, \textsc{Update}, \textsc{Delete}).

A memory query sends $q \in \mathcal{Q}$ to the retrieval server, which returns matching subtrees with path annotations:
\begin{equation}
\texttt{retrieve}(q, M_t) = \bigl\{(v, \text{path}(r \leadsto v), \text{subtree}(v)) \;\big|\; v \in V_t,\; \text{match}(q, v)\bigr\}.
\end{equation}
A memory update modifies the tree via $f_{\text{update}}: \mathcal{M} \times \mathcal{T} \to \mathcal{M}$.

At turn $t$, the agent observes input $x_t$ and memory state $\mathcal{M}_t$, and produces a response through a trajectory of interleaved language and memory actions:
$
\tau_t = (a_{t,1}, a_{t,2}, \ldots, a_{t,L}) \sim \pi_\theta(\cdot \mid x_t, \mathcal{M}_t).
$

\subsection{Verifiable Rewards}

The total reward decomposes into three verifiable components:
\begin{equation}
R(\tau_t) = \alpha \cdot R_{\text{syntax}}(\tau_t) + \beta \cdot R_{\text{utility}}(\tau_t) + \gamma \cdot R_{\text{efficiency}}(\tau_t),
\end{equation}


$R_{\text{syntax}}(\tau_t) = \mathbf{1}[\tau_t \text{ uses well-formed } \texttt{<memory>}\text{, } \texttt{<think>}\text{, } \texttt{<answer>} \text{ tags}]$ verifies format correctness.

$R_{\text{utility}}(\tau_t) = \text{BERTScore-F1}(y_t, y_t^*)$ measures answer 
correctness against the ground truth $y_t^*$ via contextual embedding similarity, 
providing a soft, model-based alternative to exact-match metrics suitable for both 
factoid and open-ended tasks.

$R_{\text{efficiency}}(\tau_t) = -\lambda_{\text{tok}} \cdot |\text{tokens}_{\text{retrieved}}| / C - \lambda_{\text{act}} \cdot |\text{mem\_actions}|$ penalizes over-retrieval, incentivizing targeted subtree access.



\subsection{Policy Optimization}

The \emph{policy} $\pi_\theta$ is a conditional distribution over action trajectories, parameterized by the LLM weights $\theta$:
\begin{equation}
\pi_\theta(\tau_t \mid x_t, M_t) = \prod_{\ell=1}^{L_t} \pi_\theta(a_{t,\ell} \mid x_t, M_t, a_{t,<\ell}),
\end{equation}
where $a_{t,<\ell} = (a_{t,1}, \ldots, a_{t,\ell-1})$ denotes the actions taken so far within the current turn. The policy factorizes autoregressively at the token level within each action. We decompose the system into two RL-trained agents.

\paragraph{Memory Manager ($\pi_{\theta_M}$).}
After each dialogue turn, the Memory Manager observes the current state $s = (x_t, M_t^{\text{old}})$ and selects an operation to maintain the memory bank:
\begin{equation}
(o, m') \sim \pi_{\theta_M}(\cdot \mid x_t, M_t^{\text{old}}), \quad o \in \{\textsc{Add},\; \textsc{Update},\; \textsc{Delete},\; \textsc{Noop}\}.
\end{equation}
The updated memory $M_t^{\text{new}} = f_{\text{update}}(M_t^{\text{old}}, o, m')$ is passed to Stage~2. The manager is rewarded based on whether its edits enable the Answer Agent to respond correctly downstream, learning to consolidate related facts (e.g., merging two pet adoption events) rather than overwrite or fragment them.

\paragraph{Answer Agent ($\pi_{\theta_A}$).}
Given a user question $q$, the Answer Agent retrieves candidate memories $\mathcal{C} = \{m_1, \ldots, m_n\}$ via RAG, then applies a memory distillation policy to filter them:
\begin{equation}
\hat{\mathcal{C}} = \text{Distill}_{\pi_{\theta_A}}(\mathcal{C}, q), \qquad y = \pi_{\theta_A}(\cdot \mid q, \hat{\mathcal{C}}).
\end{equation}
Distillation surfaces only the entries most relevant to $q$, reducing noise and improving factual accuracy. The agent is rewarded for answer correctness.


 Both agents are trained with outcome-driven reinforcement learning, requiring only downstream QA correctness as supervision.
 For a fixed context $(x_t, \mathcal{M}_t)$, we sample a group of $K$ trajectories and assign each a scalar reward $r(\tau_t^{(k)})$. Rather than using these rewards as absolute training signals, we optimize the policy using Group Relative Policy Optimization (GRPO), which computes a group-relative advantage for each trajectory:
\begin{equation}
\hat{A}^{(k)} = \frac{r(\tau_t^{(k)}) - \mu}{\sigma}, \quad \mu = \frac{1}{K}\sum_{k=1}^{K} r(\tau_t^{(k)}), \quad \sigma = \mathrm{std}\!\left(\{r(\tau_t^{(k)})\}_{k=1}^{K}\right).
\end{equation}

% Given state $s = (x_t, M_t^{\text{old}})$, the manager samples $K$ operation trajectories $\{(o^{(k)}, m'^{(k)})\}_{k=1}^K$, each producing an updated memory $M_t^{(k)} = f_{\text{update}}(M_t^{\text{old}}, o^{(k)}, m'^{(k)})$. The reward is whether the Answer Agent can answer correctly from $M_t^{(k)}$: trajectories that produce well-consolidated memory receive positive advantage, while those that fragment or overwrite useful information receive negative advantage.


%% =====================================================================
%% SECTION 2: PROGRESS TO DATE
%% =====================================================================
\section{Progress to Date}

\subsection{System Architecture}

We have implemented both Memory-R1 (flat memory) and \textsc{Struct Memory-R1} (tree-structured memory) following the two-agent architecture described in Section~2. The full implementation, including the Memory Manager, Answer Agent, data construction pipeline, GRPO training scripts, evaluation metrics, and end-to-end inference pipeline, is publicly available at: \url{https://github.com/YifanLiu2/Struct_Memory_R1.git}.


\subsection{Data Curation}

We prepare data from two sources: the LoCoMo conversational benchmark (adapted for both flat and structured memory) and three structured domains from the Semantic XPath evaluation setup.

\paragraph{Flat LoCoMo.}
We use the LoCoMo benchmark~\citep{maharana2024evaluating}, a long multi-session dialogue dataset consisting of 10 conversations ($\sim$600 turns, $\sim$26k tokens each) annotated with QA pairs spanning single-hop, multi-hop, open-domain, and temporal reasoning. Following Memory-R1, we use GPT-4o-mini to extract facts from dialogue turns and build temporal flat memory banks.

\paragraph{Structured LoCoMo.}
We convert the same LoCoMo dialogues into tree-structured memory.
The dataset contains: sessions of speaker turns with timestamps, per-session event summaries, per-session prose summaries, and per-session observations---speaker-attributed, turn-grounded atomic claims such as \textit{``Caroline is researching adoption agencies''} (D2:8). We use the following schema:

\dirtree{%
.1 Conversation.
.2 Speaker (name).
.3 Observation (text, source\_dia\_id, session\_id).
.2 Session (session\_id, datetime).
.3 Summary (text).
.3 Event (speaker, description).
.3 Turn (dia\_id, speaker, text, image\_caption).
}

\texttt{Speaker} nodes aggregate persistent observations about each participant across all sessions. \texttt{Session} nodes preserve the raw conversational record alongside two levels of derived structure: dense prose \texttt{Summary} nodes, which enable efficient session-level semantic filtering, and sparse \texttt{Event} nodes capturing key milestones. The dataset's gold observations map directly to \texttt{Observation} nodes, providing supervision for memory construction without requiring LLM-based extraction.

This schema supports two natural query patterns. \textit{Profile queries} retrieve persistent facts about a speaker directly from their observation nodes. \textit{Temporal queries} use \texttt{Summary} nodes as a predicate to narrow to the relevant session before retrieving specific turns within it.

\paragraph{Semantic XPath Domains (OOD).}
Following the evaluation setup in Semantic XPath~\citep{semanticxpath2026}, we curated three structured memory domains as out-of-distribution (OOD) test sets:
\begin{itemize}[nosep]
    \item \textbf{Travel Itinerary}: Schema \texttt{Itinerary} $\to$ \texttt{Version} $\to$ \texttt{Day} $\to$ \texttt{POI} (20 QA pairs).
    \item \textbf{To-Do List}: Schema \texttt{TodoList} $\to$ \texttt{Category} $\to$ \texttt{Project} $\to$ \texttt{Task} (20 QA pairs).
    \item \textbf{Meal Kit}: Schema \texttt{MealPlan} $\to$ \texttt{Day} $\to$ \texttt{Meal} $\to$ \texttt{Option} (20 QA pairs).
\end{itemize}
These domains generalization to unseen tree schemas and question types not present in the LoCoMo training data.



\subsection{Training}
To validate the RL training infrastructure before adapting it for memory management in \textsc{Struct Memory-R1}, we currently reproduced the training process of the foundation works in this paradigm, Search-R1 ~\citep{jin2025search} on a single consumer GPU ($\geq$24\,GB VRAM). This required reducing the model from \textsc{Qwen2.5-7B} to \textsc{Qwen2.5-0.5B}, lowering the GRPO group size from $K{=}5$ to $K{=}2$, enabling full FSDP CPU offloading, and training on a 100K-document Wikipedia subset (vs.\ 21M) indexed with FAISS. Table~\ref{tab:config} summarizes the configuration differences. We provided all the codes in the github repo listed above.

\begin{table}[h]
\centering
\caption{Search-R1 configuration: original 8-GPU vs.\ our single-GPU adaptation.}
\label{tab:config}
\small
\begin{tabular}{lcc}
\toprule
\textbf{Parameter} & \textbf{Original (8$\times$GPU)} & \textbf{Ours (1$\times$GPU)} \\
\midrule
Model & Qwen2.5-7B & Qwen2.5-0.5B \\
Train batch size & 512 & 8 \\
GRPO group size $K$ & 5 & 2 \\
Max prompt length & 4096 & 2048 \\
Max response length & 500 & 256 \\
vLLM GPU util. & 0.6 & 0.3 \\
FSDP CPU offload & partial & full \\
Corpus size & 21M docs & 100K docs \\
\bottomrule
\end{tabular}
\end{table}
The result shows a promising directions. This single-GPU run confirmed that the GRPO training loop, vLLM rollout generation, reward computation, and retrieval server all function correctly end-to-end. While the corpus size is much smaller than Search-R1, it is close to the training size in the Memory-R1 variant.x 

For the \textsc{Struct Memory-R1} adaptation, we will using the same config but replace the search retriever with the memory states, swap the single-agent multi-turn loop for two decoupled one-shot agents as we described above, and construct training data from LoCoMo dialogues instead of NQ question-answer pairs. The same veRL infrastructure and GRPO objective are reused throughout.


\subsection{Experiment Setup}

For LoCoMo, we split the QA pairs into train / validation / test following Memory-R1's 1:1:8 ratio. Memory-R1 (Flat) is trained on flat LoCoMo, and \textsc{Struct Memory-R1} is trained on structured LoCoMo. The three Semantic XPath domains are used only for OOD evaluation. For the base LLM, we use Qwen2.5-0.5B, the same model used in our Search-R1 preliminary verification (Section~3.2), with the same single-GPU GRPO configuration adapted for memory operations.

\subsection{Baselines and Methods}

\paragraph{Non-RL Methods.}
\textbf{Zero-Shot (no memory)}: The base LLM answers directly without access to any external memory. This serves as a lower bound.

\textbf{In-Context Memory}: The entire memory state is serialized and appended to the LLM prompt. Provides maximal information but scales poorly due to the ``lost in the middle'' phenomenon~\citep{liu2024lost}.

\textbf{Flat RAG}: Standard retrieval-augmented generation with embedding-based top-$k$ retrieval over the flat memory bank, without RL training.

\textbf{Semantic XPath}~\citep{semanticxpath2026}: The LLM generates an XPath-style query over the structured memory tree, executed by the Semantic XPath engine. We use results reported in the Semantic XPath paper as reference.

\paragraph{RL Methods.} \textbf{Memory-R1 (Flat)}~\citep{yan2025memory}: The Memory-R1 two-agent architecture operating on flat memory with ADD/UPDATE/DELETE/NOOP operations, trained with GRPO.

\textbf{\textsc{Struct Memory-R1} (ours)}: Our proposed extension of Memory-R1 to tree-structured memory, where ADD operations require a \texttt{parent\_id} specifying placement in the hierarchy, trained with GRPO.


\subsection{Signs of Life: Sample Trajectories}

To demonstrate that our framework shows signs of life, we compare single trajectories from three methods on representative queries.

% \paragraph{Trajectory 1: Zero-Shot (No Memory).}
\begin{tcolorbox}[aclbox, width=\linewidth]
\small
\textbf{Trajectory 1: Zero-Shot (No Memory).}\\begin{equation}2pt]
\textbf{Q:} How many pets does Alice have?\\begin{equation}2pt]
\textbf{LLM:} Based on my knowledge, I don't have specific information about Alice's pets. A common answer might be one or two.\\begin{equation}2pt]
\textbf{Output:} \texttt{**Answer:** 1}\\begin{equation}2pt]
\textbf{Result:} \textcolor{red}{Incorrect} (gold: 2 dogs). The model hallucinates without access to the conversation history.
\end{tcolorbox}

% \paragraph{Trajectory 2: Memory-R1 (Flat).}
\begin{tcolorbox}[aclbox, width=\linewidth]
\small
\textbf{Trajectory 2: Memory-R1 (Flat).}
\textbf{Stage 1 -- Memory Manager processes dialogue turn:}\\
\textit{Input:} Facts: [``Alice adopted two dogs named Buddy and Scout'']\\
\textit{Current bank:} [\{id:0, text:``Alice works as a software engineer''\}, \{id:1, text:``Alice likes cheese pizza''\}]\\
\textit{MM output:} \texttt{\{``memory'': [\{``id'':``0'', ``event'':``NONE''\}, \{``id'':``1'', ``event'':``NONE''\}, \{``id'':``2'', ``text'':``Alice adopted two dogs named Buddy and Scout'', ``event'':``ADD''\}]\}}\\[4pt]
\textbf{Stage 2 -- Answer Agent:}\\
\textit{Retrieved memories (top-3):} ``Alice adopted two dogs named Buddy and Scout'' | ``Alice works as a software engineer'' | ``Alice likes cheese pizza''\\
\textit{AA output:} Selected: Memory about dog adoption.\\
\texttt{**Answer:** 2 dogs}\\[2pt]
\textbf{Result:} \textcolor{green!50!black}{Correct}. The flat manager correctly ADDs the new fact; the Answer Agent distills the relevant memory.
\end{tcolorbox}

% \paragraph{Trajectory 3: \textsc{Struct Memory-R1} (Tree).}
\begin{tcolorbox}[aclbox, width=\linewidth]
\small
\textbf{Trajectory 3: \textsc{Struct Memory-R1} (Tree).}
\textbf{Stage 1 -- Tree Memory Manager processes dialogue turn:}\\
\textit{Current tree:}\\
\texttt{~~[Dialogue] (id=root)}\\
\texttt{~~~~[Session] (id=s1) text: Morning chat} | 
\texttt{[Memory] (id=m1) text: Alice is a software engineer}\\
\texttt{~~~~[Session] (id=s2) text: Evening catchup} | 
\texttt{[Memory] (id=m3) text: Bob is a teacher}\\[2pt]
\textit{Facts:} [``Alice adopted two dogs named Buddy and Scout'']\\
\textit{MM output:} \texttt{\{``memory'': [\{``id'':``m4'', ``parent\_id'':``s1'', ``node\_type'':``Memory'', ``text'':``Alice adopted two dogs: Buddy and Scout'', ``event'':``ADD''\}]\}}\\[4pt]
\textbf{Key:} The manager places the new node under \texttt{s1} (Alice's session), not \texttt{s2} (Bob's session) -- a structural decision unique to \textsc{Struct Memory-R1}.\\[4pt]
\textbf{Stage 2 -- Answer Agent:}\\
\textit{Retrieved memories (with paths):} [Path: /Dialogue/Session(s1)/Memory(m4)] ``Alice adopted two dogs: Buddy and Scout''\\
\texttt{**Answer:** 2 dogs}\\[2pt]
\textbf{Result:} \textcolor{green!50!black}{Correct}. The tree path context helps disambiguate whose pets are being asked about.
\end{tcolorbox}

\subsection{Preliminary Results}

Table~\ref{tab:preliminary} results across all OOD datasets. Results for the RL methods are pending completion of GRPO training.

\begin{table}[h]
\centering
\caption{Preliminary results across evaluation datasets use Pass Rate. ``/'' indicates results pending.}
\label{tab:preliminary}
\small
\begin{tabular}{lccccc}
\toprule
\textbf{Method} & \textbf{Flat LoCoMo} & \textbf{Struct LoCoMo} & \textbf{Itinerary} & \textbf{To-Do} & \textbf{Meal Kit} \\
\midrule
\multicolumn{6}{l}{\textit{Non-RL Methods}} \\
Zero-Shot (no memory) & / & / & 15.0\% & / & / \\
In-Context Memory & / & / & 70.0\% & 70.0\% & 65.0\% \\
Flat RAG & / & / & 30.0\% & / & / \\
Semantic XPath & -- & / & 85.0\% & 80.0\% & 85.0\% \\
\midrule
\multicolumn{6}{l}{\textit{RL Methods}} \\
Memory-R1 (Flat) & / & -- & / & / & / \\
\textsc{Struct Memory-R1} (ours) & -- & / & / & / & / \\
\bottomrule
\end{tabular}
\end{table}

The Semantic XPath baselines on the OOD domains achieve 80--85\% pass rates with significantly fewer tokens than In-Context Memory, establishing strong non-RL reference points. \textsc{Struct Memory-R1} targets the structured LoCoMo and OOD domains where tree-aware memory management can leverage hierarchical relationships.

%% =====================================================================
%% SECTION 3: THE PLAN MOVING FORWARD
%% =====================================================================
\section{The Plan Moving Forward}

Our immediate next step is to complete GRPO training for both Memory-R1 (flat) and \textsc{Struct Memory-R1} (tree) using Qwen2.5-0.5B as the base model. Specifically, we plan to:
\begin{enumerate}[nosep]
    \item Run the two-stage training pipeline (Memory Manager $\to$ Answer Agent) on the LoCoMo training split for both flat and structured memory configurations.
    \item Evaluate on all five datasets (flat LoCoMo, structured LoCoMo, and the three OOD domains) using F1, BLEU-1, Sub-EM, and LLM-as-a-Judge metrics.
    \item Analyze the results to understand the effect of structured vs.\ flat memory on different question types (single-hop, multi-hop, temporal, open-domain) and on in-distribution vs.\ OOD generalization.
\end{enumerate}

\bibliographystyle{plainnat}
\bibliography{reference}

\end{document}